\section{Related Work}\label{sec:related}

\subsection{On-Device Deep Learning Benchmarking}
Early mobile AI measurements established the extreme performance diversity of smartphone accelerators. Ignatov et al.~\cite{ignatov2018ai,ignatov2019ai} introduced AI Benchmark across Android handsets, exposing wide latency disparities across CPU, GPU, and DSP/NPU backends. MLPerf Mobile~\cite{reddi2020mlperf,reddi2022mlperfmobile} standardized open-source accuracy and latency rules for edge vision and NLP models. Measuring in-the-wild deployment, Xu et al.~\cite{xu2019first} mined $16{,}500$ top Android apps and identified $211$ deep learning apps, noting rapid growth in embedded vision models such as MobileNet and SSD. Almeida et al.~\cite{almeida2021smart} compared device tiers and observed low-end and mid-tier phones running on average $3.4\times$ and $1.51\times$ slower than flagship silicon, showing non-linear FLOP-to-latency behavior driven by thermal throttling, scheduler policies, and excessive CPU thread contention. Recent benchmarks extend these measurements to generative mobile workloads, including MobileAIBench~\cite{murthy2024mobileaibench} and lm-Meter~\cite{wang2025lmmeter}, profiling end-to-end latency and memory dynamics of on-device language and multimodal models inside native C++/Android harnesses.

\subsection{Cross-Platform Mobile Framework Performance}
Bi{\o}rn-Hansen et al.~\cite{biornhansen2019survey} classified cross-platform architectures into WebView hybrids, JavaScript bridges, and natively compiled toolkits. Earlier comparisons by Willocx et al.~\cite{willocx2016comparing} and Del{\'\i}a et al.~\cite{delia2015multi} measured cold-launch delays and synthetic compute overheads across early hybrid stacks, while Ciman and Gaggi~\cite{ciman2017empirical} and Horn et al.~\cite{horn2023native} quantified persistent frame and energy costs relative to native Android. Focusing on modern AOT-compiled UI toolkits, Bi{\o}rn-Hansen et al.~\cite{biornhansen2020empirical} showed that Flutter approaches native UI frame responsiveness at the cost of higher memory consumption from the Dart VM and its Skia rendering layer; subsequent Flutter releases have migrated rendering from Skia to Impeller~\cite{flutterimpeller2025} (Vulkan on Android API~29+, OpenGL ES fallback on API~$\le 28$). Nanavati et al.~\cite{nanavati2024critical} further examined Flutter runtime profiling and memory tuning for UI-bound applications.

\subsection{Mobile Runtimes, Quantization, and Benchmark Rigor}
Inference efficiency depends strongly on native library implementation and numerical precision. Zhang et al.~\cite{zhang2022comprehensive} benchmarked six mobile engines across 15 models and 10 devices, finding that engine choice swings average latency by $7.4\times$ overall ($8.7\times$ CPU, $6.0\times$ GPU, and up to $54.3\times$ worst-case) while INT8 quantization delivers $0.8\text{--}3.0\times$ speedups; Luo et al.~\cite{luo2020comparison} similarly highlighted operator-coverage fragility across delegates. Jacob et al.~\cite{jacob2018quantization} formalized integer-only quantization to cut parameter footprints $4\times$ and accelerate SIMD hardware, yet delegation gains remain bounded by shader initialization and copy costs. Methodologically, Georges et al.~\cite{georges2007statistically} and Barrett et al.~\cite{barrett2017virtual} showed that managed-runtime benchmarks require separating warm-up transients from steady-state measurements, motivating Cliff's~$\delta$~\cite{cliff1993dominance} non-parametric dominance effects, bootstrap intervals, and Holm--Bonferroni family-wise control~\cite{holm1979simple}.

\subsection{Combined Research Gap}\label{subsec:gap}
Prior literature remains bifurcated across two silos. Machine learning benchmarks~\cite{ignatov2019ai,reddi2022mlperfmobile,zhang2022comprehensive,wang2025lmmeter} isolate C++ engines from host UI application frameworks, whereas cross-platform studies~\cite{biornhansen2020empirical,horn2023native,nanavati2024critical} test scrolling lists or synthetic loops without neural networks, delegates, or quantization. Consequently, no prior work isolates how a modern Flutter stack alters every pipeline stage---pure-Dart preprocessing, FFI marshalling, upstream dependency drift, and heap vs.\ \texttt{mmap} model loading---relative to idiomatic native Android under identical pinned models. Our protocol addresses this combined gap on entry-tier hardware using warm-up isolation, non-parametric equivalence tests, and session-level aggregation.
